\section{Synthesis: what the atlas hands to theory}
\label{sec:synthesis}

\manuscriptfigure{figures/atlas_map.pdf}{%
  \textbf{Verdict map.} One row per family: block colour, the outcome returned
  by the decision procedure, the evidence grade, and the six-component evidence
  vector. Outcomes are always written as text; colour is redundant with the
  label rather than carrying it.%
}{fig:atlas-map}

Figure~\ref{fig:atlas-map} is the release in one view. Of
\AtlasTotalFamilies{} families, the procedure returned a spread rather than a
uniform verdict, and the spread is the point: a method that returns
``supported'' everywhere has not been tested.

\subsection{The calibration result}

The single most informative entry in the map is which family passed. Exactly
\AtlasCountPhaseTransitionSupported{} of \AtlasTotalFamilies{} families
received the verdict \emph{transition supported}, and it is the Gardner
capacity --- the anchor whose answer was known before the experiment was run
(\S\ref{sec:anchors}). Applied to the fourteen settings whose answers are not
known, the same procedure certified nothing.

We take that as the central evidence that the gates are calibrated rather than
merely strict. A procedure that never fires is useless; one that fires
everywhere is worse. This one fired once, on the case placed there to make it
fire, and the boundary it reported,
$\alpha=\AtlasPerceptronCapacityPeakControl{}$, sits at the grid point nearest
the exact $\alpha_c=2$.

\subsection{Three transferable findings}

\paragraph{1. The control parameter is a hypothesis, and in the
best-resolved family it fails.}
The strongest critical response in the release --- an interior susceptibility
peak at every size, with height growing as
$N^{\AtlasSpecializationSharpening{}}$ --- comes with peak locations that drift
systematically in $\alpha=n/d$. The fitted path exponent
$\gamma=\AtlasSpecializationGamma{}$, with $95\%$ interval
\AtlasSpecializationGammaCI{}, excludes the proportional regime $\gamma=1$.
The immediate consequence is methodological rather than physical: a study that
sweeps $\alpha$ at a single size is not holding the boundary fixed, and two
such studies at different sizes can disagree without either being wrong.
This is a quantitative, falsifiable target --- if a theory predicts a fixed
$\alpha_c$ for this setting, it is in tension with \eqref{eq:gamma-spec}.

\paragraph{2. Diverging response and locatable boundary are different
achievements, and the second is much harder.}
Several families produce a susceptibility that grows convincingly with size and
still fail to yield a boundary location: the peak marches out of the sampled
window, or the drift extrapolation and the collapse name places that differ by
more than the consistency tolerance. Reporting only the first --- the usual
practice --- would have produced several confident transition claims here that
the location analysis does not support. We suggest that a susceptibility plot
without an accompanying drift fit should be read as evidence of criticality and
not as evidence of a critical point.

\paragraph{3. Two design choices can manufacture or destroy a transition.}
Both were found the hard way in this work. A control grid too coarse to hold an
interior maximum yields a structurally censored result regardless of the
physics; and a fixed epoch budget makes wider models appear \emph{less} ordered,
because they are less optimised, which inverts the sign of the finite-size
trend. Any finite-size claim about a learning system should therefore report
grid density and the optimisation stopping rule alongside the seed count. We
recommend training to a convergence criterion, and recording per run whether it
was met, so that optimisation shortfall is a measured quantity rather than a
hidden covariate.

\subsection{The ledger}

\begin{table}[tbp]
\centering\small
\begin{tabularx}{\linewidth}{@{}p{34mm}p{20mm}Y@{}}
\toprule
Outcome & Families & Reading \\
\midrule
transition supported & \AtlasCountPhaseTransitionSupported{} & two consistent
finite-size signatures and a located boundary \\
critical response, boundary unresolved &
\AtlasCountCriticalResponseBoundaryUnresolved{} & the response diverges; the
window does not contain the boundary \\
candidate boundary & \AtlasCountCandidateBoundary{} & one admissible signature;
a second is needed \\
response shift (off-grid) &
\AtlasCountResponseShiftBoundaryOutsideGrid{} & the order parameter moves;
most sizes censored \\
no resolved response & \AtlasCountNoResolvedResponse{} & the order parameter
does not separate from its uncertainty \\
\bottomrule
\end{tabularx}
\caption{Outcome ledger for the release. Rows below the first are reported with
the same prominence as the first; a map whose empty cells are hidden is not a
map.}
\label{tab:ledger}
\end{table}

\subsection{What we would do next, in order}

\begin{enumerate}
  \item \textbf{Twelve seeds on the four best-resolved families.} Every family
    here is capped at grade~C purely by the replication tier. This is the
    cheapest available upgrade and changes nothing else.
  \item \textbf{Re-run the drifting families on the rescaled control
    $n/d^{\gamma}$.} If the peaks return to the interior of the window, several
    unresolved boundaries become locatable in a single sweep.
  \item \textbf{A frozen prospective replication.} No result here has been
    confirmed at fresh seeds after the analysis was fixed, which is the one
    gate that no amount of re-analysis can substitute for.
  \item \textbf{One natural-data endpoint.} External validity is zero across
    the release by construction. A single family carried through to a natural
    corpus would establish whether any of these boundaries survive contact with
    real data, which is the question the whole programme is ultimately for.
\end{enumerate}
